% =============================================================
% Chapitre 3 — Analyse des besoins et conception de l'architecture
% cible (~13 pages). Contenu dérivé du HLD/LLD réels du projet.
% =============================================================
\chapter{Analyse des besoins et conception de l'architecture cible}
\label{ch:conception}

\begin{attentedoc}
Ce chapitre décrit la \textbf{cible} de conception telle que fixée dans le document
d'architecture (HLD) du projet. Il doit être relu et amendé à la lumière de
\texttt{06\_ECARTS\_IMPLEMENTATION.md} avant validation finale, pour signaler clairement
tout élément qui aurait changé entre la conception et la mise en œuvre --- celle-ci étant
traitée au chapitre~\ref{ch:mise-en-oeuvre}.
\end{attentedoc}

\section{Cadrage et hypothèses structurantes}

Le cadrage retenu au démarrage du projet repose sur trois hypothèses assumées, énoncées
explicitement plutôt que découvertes après coup :

\begin{enumerate}
  \item trois environnements applicatifs (développement, recette, production) correspondant
        chacun à un projet GCP distinct, dans une seule région, complétés par un projet
        d'amorçage ne contenant aucun service applicatif (état Terraform, fédération
        d'identité, registre d'images) ;
  \item la montée en charge n'est pas un objectif du projet --- les limites de ce choix sont
        documentées plutôt que contournées par une architecture surdimensionnée ;
  \item la simplicité est traitée comme un critère d'architecture à part entière : tout
        composant proposé doit justifier son existence, et tout composant non justifié est
        retiré.
\end{enumerate}

\section{Notions clés}

Avant de détailler les besoins puis l'architecture, il est utile de fixer le sens de quelques
termes employés tout au long de ce chapitre, pour un lecteur qui ne serait pas nécessairement
spécialiste du cloud ou de la cybersécurité.

\textbf{Zero Trust} désigne un modèle de sécurité dans lequel aucune entité ne bénéficie d'une
confiance implicite du seul fait de sa position réseau : chaque accès est vérifié explicitement,
quelle que soit son origine (cf. chapitre~\ref{ch:etat-art}).

\textbf{DevSecOps} désigne l'intégration des contrôles de sécurité directement dans le pipeline
de livraison logicielle, sous forme de portes automatisées, plutôt que sous forme de revues
manuelles réalisées après coup.

\textbf{Infrastructure as Code (IaC)} désigne la description de l'infrastructure sous forme de
code déclaratif versionné, permettant de reproduire, revoir et faire évoluer un environnement
de la même manière qu'un programme.

\textbf{Micro-segmentation par identité} désigne, dans le contexte de ce projet, une approche
où le contrôle des communications entre composants ne repose pas sur un découpage réseau
(VLAN, sous-réseaux), mais sur l'identité vérifiée de chaque composant --- un compte de service
distinct par élément de l'architecture, avec des droits strictement délimités.

\textbf{Embedding sémantique} désigne la transformation d'un texte (ici, un événement de
sécurité) en un vecteur numérique, de telle sorte que deux événements proches en sens le
soient également dans l'espace vectoriel --- ce qui permet de rapprocher un événement d'une
technique d'attaque connue même sans correspondance exacte de motif.

\section{Expression des besoins}

\begin{table}[htbp]
\centering
\caption{Exigences non fonctionnelles retenues pour le socle}
\label{tab:exigences-nf}
\begin{tabular}{@{}P{3cm}P{6cm}P{5.5cm}@{}}
\toprule
\textbf{Exigence} & \textbf{Cible} & \textbf{Remarque} \\
\midrule
Disponibilité & Objectifs de niveau de service instrumentés sur les services critiques &
Mesuré en continu, pas de simple \og best effort \fg{} \\
Scalabilité & Volontairement limitée (cadrage projet) & Limites connues et documentées \\
Sécurité & Conformité Zero Trust (NIST SP~800-207) ; DevSecOps avec portes bloquantes &
Cœur du projet \\
Coût & Maîtrisé (paliers gratuits / faible coût mensuel) & Poste dominant : ingestion de
journaux \\
Auditabilité & Chaque décision automatisée journalisée (hachage, score, version) &
Prérequis de la valeur \og SIEM \fg{} \\
\bottomrule
\end{tabular}
\end{table}

Sur le plan fonctionnel, l'architecture doit répondre à quatre besoins directement issus de
la problématique du chapitre~\ref{ch:contexte} : centraliser la gestion des identités et des
autorisations plutôt que de les dupliquer par application ; contrôler les communications
entre composants au niveau de l'identité plutôt qu'au niveau du réseau ; bloquer
automatiquement toute livraison logicielle ne respectant pas un ensemble de règles de
sécurité définies ; et rattacher les événements de sécurité observés aux techniques
d'attaque connues, y compris lorsqu'aucune règle statique ne les couvre.

\section{Vue d'ensemble de l'architecture cible}

L'architecture cible est organisée en sept couches, complétées par un plan de livraison
transversal (figure~\ref{fig:archi-globale}) :

\begin{table}[htbp]
\centering
\caption{Couches de l'architecture cible}
\label{tab:couches-archi}
\begin{tabular}{@{}P{1.2cm}P{3.3cm}P{9cm}@{}}
\toprule
\textbf{Couche} & \textbf{Nom} & \textbf{Rôle} \\
\midrule
L1 & Edge & Point d'entrée unique : résolution de nom, répartition de charge, pare-feu
applicatif et limitation de débit \\
L2 & Identité \& sécurité & Plan transversal : comptes de service, gestion des secrets,
chiffrement, autorisation applicative \\
L3 & Workloads & Services exécutés à la demande (sans serveur à administrer) \\
L4 & Réseau & Réseau privé, accès privé aux services de données, sorties contrôlées \\
L5 & Données & Base de données applicative et entrepôt de détection (SIEM) \\
L6 & ML / enrichissement sémantique & Encodage vectoriel et rattachement aux techniques
d'attaque connues \\
L7 & Observabilité & Plan transversal : journalisation, alerte, tableaux de bord \\
\bottomrule
\end{tabular}
\end{table}

\begin{attentedoc}
Insérer ici le schéma d'architecture détaillé (figure~\ref{fig:archi-globale}), exporté
depuis le document source du projet.
\end{attentedoc}

% PLACEHOLDER — à remplacer par le schéma réel dès disponible :
% \begin{figure}[htbp]
%   \centering
%   \includegraphics[width=\linewidth]{figures/archi_globale.pdf}
%   \caption{Architecture cible (7 couches + plan de livraison transversal)}
%   \label{fig:archi-globale}
% \end{figure}
\begin{figure}[htbp]
  \centering
  \fbox{\begin{minipage}[c][6cm][c]{0.9\linewidth}
    \centering
    \textcolor{red!60!black}{\textbf{[Schéma en attente --- voir encadré ci-dessus]}}
  \end{minipage}}
  \caption{Architecture cible --- placeholder, à remplacer par le schéma source}
  \label{fig:archi-globale}
\end{figure}

\section{Application des principes Zero Trust}

\begin{table}[htbp]
\centering
\caption{Traduction des principes Zero Trust dans l'architecture MENAL}
\label{tab:zt-principes}
\begin{tabular}{@{}P{4.2cm}P{9cm}@{}}
\toprule
\textbf{Principe (NIST 800-207)} & \textbf{Traduction dans l'architecture} \\
\midrule
Ne jamais faire confiance au réseau & Aucune ressource de données n'est joignable
publiquement ; être \og à l'intérieur \fg{} ne donne aucun droit \\
Vérifier explicitement à chaque requête & Chaque appel entre services porte une identité
vérifiée ; aucun appel anonyme entre composants internes \\
Moindre privilège & Un compte de service distinct par composant, avec des droits limités
au strict nécessaire \\
Présomption de compromission & Toute action est journalisée ; un composant d'analyse ne
peut jamais modifier les preuves qu'il traite \\
Micro-segmentation & Séparation par identité plutôt que par segment réseau \\
Surveillance continue & Journalisation et alerte sur l'ensemble des couches, y compris la
couche d'enrichissement \\
\bottomrule
\end{tabular}
\end{table}

\section{Conception de la stratégie d'identité, de réseau et du pipeline}

Un compte de service distinct est attribué à chaque composant, sans clé d'accès exportée :
l'authentification du pipeline de livraison auprès de la plateforme cloud repose sur une
fédération d'identité plutôt que sur un secret statique stocké dans le système
d'intégration continue --- un choix directement lié au constat du
chapitre~\ref{ch:etat-art} sur les risques de la chaîne de livraison.

\begin{alertbox}
Un principe de conception est appliqué strictement au niveau des droits d'accès : le
composant chargé de l'enrichissement sémantique (couche~L6) peut \emph{lire} les événements
de détection, mais ne dispose d'\emph{aucun} droit d'écriture sur les journaux bruts ni sur
les détections elles-mêmes --- seulement sur la table dédiée à ses propres résultats. Un
moteur de détection ne doit jamais pouvoir modifier les preuves qu'il analyse.
\end{alertbox}

\todo{Détailler ici la matrice complète des comptes de service et de leurs rôles --- la
version complète est reportée en annexe~D pour respecter la contrainte de concision du
corps du rapport ; n'en garder ici qu'un tableau réduit aux 4--5 comptes les plus
illustratifs.}

\section{Conception de la couche d'enrichissement sémantique}

La couche d'enrichissement (L6) répond au besoin identifié au chapitre~\ref{ch:etat-art} :
rattacher un événement de sécurité à une technique d'attaque connue même en l'absence de
règle statique correspondante. Le choix retenu privilégie un modèle d'embedding
\emph{non génératif}, dont la sortie est un vecteur déterministe comparé par similarité aux
techniques du référentiel MITRE ATT\&CK, plutôt qu'un modèle de langage génératif dont la
sortie serait non déterministe et potentiellement influençable par le contenu même des
journaux analysés --- une propriété jugée incompatible avec une chaîne de preuve.

Ce composant est positionné pour n'être joignable que depuis l'intérieur du réseau privé, et
son résultat (technique rattachée, score de similarité, version du modèle) est systématiquement
journalisé, ce qui permet d'auditer une décision d'enrichissement a posteriori.

\section{Registre des décisions d'architecture}

Chaque choix de conception a été examiné puis classé selon un statut explicite plutôt
qu'accumulé sans justification --- une pratique directement héritée de la revue
d'architecture menée en amont du projet.

\begin{table}[htbp]
\centering
\caption{Extrait du registre de décisions d'architecture (ADR)}
\label{tab:adr-extrait}
\begin{tabular}{@{}P{1.5cm}P{5cm}P{1.8cm}P{5.5cm}@{}}
\toprule
\textbf{Réf.} & \textbf{Décision} & \textbf{Statut} & \textbf{Justification} \\
\midrule
ADR-01 & Exécution serverless comme unique plateforme d'exécution & Validé & Pas de système
d'exploitation à durcir, montée à zéro possible, identité native par révision \\
ADR-04 & Entrepôt de données comme moteur du SIEM (plutôt qu'une solution commerciale) &
Validé & Le SIEM est construit, pas acheté --- c'est la valeur pédagogique du projet \\
ADR-05 & Modèle d'embedding non génératif pour l'enrichissement & Validé & Sortie
déterministe = preuve auditable ; pas de surface d'injection \\
ADR-08 & Fédération d'identité pour l'authentification du pipeline & Ajouté & Supprime la
classe de risque \og clé statique exportée \fg{} \\
ADR-12 & Infrastructure entièrement décrite en code, revue avant application & Validé
(renforcé) & Répond directement à un besoin de reproductibilité \\
\bottomrule
\end{tabular}
\end{table}

\begin{attentedoc}
Compléter ce tableau avec l'ensemble des décisions du registre projet (annexe possible si la
liste complète dépasse une page) --- vérifier qu'aucune décision listée ici n'a été révisée
depuis la rédaction du document source.
\end{attentedoc}

Trois décisions structurantes gagnent à être présentées sous forme d'alternative comparée
plutôt que de simple affirmation, pour rendre l'arbitrage lisible :

\begin{table}[htbp]
\centering
\caption{Alternatives comparées pour trois décisions structurantes}
\label{tab:alternatives-comparees}
\begin{tabular}{@{}P{3.3cm}P{4.7cm}P{4.7cm}@{}}
\toprule
\textbf{Décision} & \textbf{Option retenue} & \textbf{Option écartée} \\
\midrule
Plateforme d'exécution & Serverless (ADR-01) : pas de système à durcir, montée à zéro &
Cluster orchestré (Kubernetes managé) : flexibilité supérieure, mais un plan de contrôle
entier à opérer pour un périmètre qui n'en a pas besoin \\
Moteur du SIEM & Entrepôt de données existant, requêtes de détection en langage
déclaratif (ADR-04) & Solution SIEM commerciale ou open source dédiée : plus de
fonctionnalités prêtes à l'emploi, mais masque la valeur pédagogique --- le SIEM est ici
\emph{construit}, pas acheté \\
Enrichissement des alertes & Modèle d'embedding non génératif (ADR-05) & Modèle de langage
génératif : plus \og impressionnant \fg{} en démonstration, mais sortie non déterministe et
non auditable --- incompatible avec une chaîne de preuve \\
\bottomrule
\end{tabular}
\end{table}

\section{Éléments volontairement écartés}

Certains composants examinés puis écartés méritent d'être mentionnés, car documenter un
refus motivé est aussi révélateur d'une démarche d'architecte que documenter un choix
retenu :

\begin{itemize}
  \item un périmètre de contrôle des flux sortants à l'échelle de l'organisation, jugé
        disproportionné pour un projet mono-projet de démonstration ;
  \item un maillage de services avec chiffrement mutuel dédié, la plateforme d'exécution
        retenue fournissant déjà un chiffrement en transit et une identité par service ;
  \item un cache distribué, nécessaire uniquement pour une exigence de montée en charge
        explicitement hors objectif ;
  \item un modèle de langage génératif pour l'analyse des incidents, écarté par principe
        (cf. \S\ref{sec:zero-trust-nist} et raisonnement du \S\ref{ch:conception} ci-dessus
        sur la non-reproductibilité) plutôt que par contrainte technique.
\end{itemize}
